% Inactive modular snapshot; edit the root neurips_2026.tex instead.
\section{Introduction}
\label{sec:intro}

Parameter-efficient fine-tuning (PEFT) adapts a pretrained model through a restricted set of trainable parameters. Rank, parameter count, and task performance are natural descriptions of these restrictions, but they leave an important question open: where does the effective update lie relative to the structure of the pretrained weights? Two updates can have the same rank and magnitude while modifying different pretrained directions. Likewise, similar downstream performance need not imply similar spectral changes, as studies comparing LoRA and full fine-tuning demonstrate~\citep{hu2022lora,shuttleworth2025illusion}.

The singular value decomposition of a pretrained weight matrix provides a fixed reference for this question. Its leading and tail singular subspaces define a partition of the update into leading--leading, tail--tail, and two cross blocks. This is a geometric distinction, not a semantic assignment: a large singular value does not by itself identify stored knowledge, and a small singular value does not identify dispensable noise. The reference frame instead lets us ask what a constraint preserves, what freedom it removes, and how a relaxation interacts with the original directions.

This reference frame is already used by spectral fine-tuning methods, including singular-value tuning, singular-vector coefficient tuning, and regularized adaptation of singular factors~\citep{han2023svdiff,lingam2024svft,cao2025sorsa}. Our question is not whether pretrained spectral information can be used for adaptation, but what different uses of that information actually constrain. In particular, orthonormality of trainable factors, confinement to a pretrained subspace, and stability of the largest singular directions are different properties.

We investigate these distinctions through a hierarchy of spectral constraints. Fixed tail coefficients restrict both the available span and the directions expressible inside it. Within-tail rotations relax the second restriction without relaxing the first. Ambient diagonal scaling acts outside the singular coordinate system and can introduce leading--tail interaction. The implementation names LinLoRA, OrthoLoRA, UILinLoRA, and UIOrthoLoRA identify these constructions. We use them to isolate representational freedoms and study regularization of their interactions with the pretrained spectrum.

The analysis yields two separations. First, exact projection identities distinguish an outside-tail approximation residual from the additional error imposed by a diagonal tail core. They hold for any target perturbation, without assuming that task-relevant directions align with the tail. Second, the meaning of mixing control depends on the core being scaled. A tail-only core recovers confinement when mixing vanishes. A practical leading-plus-tail core instead permits direct leading adaptation even when its cross blocks vanish. Thus suppressing mixing need not preserve the leading weight block, and preserving that block need not preserve its rank ordering after adaptation.

We then examine a mixing-regularization intervention across five GLUE tasks. Penalizing both sides reduces cross energy from approximately 43\% to 10--27\% and increases tail allocation. Yet the leading block's share rises in four of five task means. This distinction would be hidden by reporting only aggregate off-tail leakage. Leading-subspace drift also falls in every paired task/seed comparison, accompanied by substantial update shrinkage and task-mean score decreases of $0.2$--$1.0$ points. The experiment establishes a geometry--magnitude--performance trade-off; it does not isolate an effect beyond magnitude-only regularization.

Our contributions are a characterization of the two confinement bottlenecks, unified bounds for coordinate-induced mixing with and without a leading core, and an empirical analysis reporting geometry, magnitude, and adaptation together. Broader language-understanding and generation results establish that the constructions support useful learning. The resulting perspective complements rank and parameter count with a precise account of what a spectral constraint restricts, what a relaxation changes, and what a regularizer can certify.
